\chapter{Technical Implementation} \label{sec:TechnicalImplementation}

\section{Engine And Frontend} \label{sec:EngineAndFrontend}

The Praxsmith project consists of an engine written in Rust and a frontend written
using the Svelte JavaScript framework that allows for developers and players to
experience the projects they write. The engine is compiled to WASM to allow it to
run directly on the browser, and TypeScript bindings were also generated to allow
for more confident frontend development.

There were a few motivating factors behind this split-architecture design:

\begin{itemize}
    \item Rust has many features that help greatly with the development of a programming
    language, including rigid type checking rules, pattern matching, and errors as values.
    This allowed for cleaner logic and fewer issues while developing, which greatly
    decreased production time.
    \item Centralizing in Rust provides the opportunity for integration with other tools,
    like the Unity game engine. (TODO: cite?) Bindings from Rust to other languages are
    widely available, such as the ``wasm-bindgen'' crate used for the web interface. As of
    now, there are only core and WASM crates, but it is trivial to add new ones that
    act as a thin wrapper around the core API.
    \item Rust's ecosystem is fairly developed, including the ``pest'' parsing library,
    which has strong integrations with Rust's compile-time type system.
    \item Rust is fast, which isn't crucial for a dialog system but helps in some scenarios,
    such as when using higher recursive depths for the planning system described below.
    \item The prevalence of WASM means that JavaScript bindings in particular are a
    well-documented technical challenge.
    \item Making the project accessible directly from the web allows for high ease of use,
    as research participants, developers, and players do not need to install anything to
    use Praxsmith. It also means that updates are immediately available on the user's
    end upon reloading the page.
    \item The developer was familiar both Rust and Svelte.
\end{itemize}

\section{Parsing} \label{sec:Parsing}

To define the rules of the language, the ``pest'' library was used. The ``pest''
language used to define the syntax is also directly used by the parsing system.
Upon parsing, ``pest'' returns detailed errors if something goes wrong, which
turned out to be crucial when developing scenarios, especially for those in the
process of learning the language.

The language is parsed from two files: a type script and a world script. The type
script is parsed into a series of relation types that are added to a type registry
upon world creation. The world script is parsed into a series of agent initializations
and relation declarations which are added to the world's internal representation
one by one.

This split script design allows for separation of the rules that guide characters
and those characters themselves. The rules can be referred to when writing the
characters, and concerns are naturally separated across this boundary. It also
allows for the reuse of rules within different stories. With more usage, a common
types file could be developed, including all that is needed to set up a basic
story with all manners of interaction.

\section{Internal Graph Representation} \label{sec:InternalGraphRepresentation}

Internally, the world is represented as a graph with agents and relations as nodes
and edges existing between these two types that connect agents to the relations they
are involved in and vice versa.

Agents are maintained as a simple hashmap with their ID as the key, which works
well because agents are initialized once upon world load and never added or removed
during runtime.

Relations need a slightly different solution, as they can be dynamically added and
removed. For this, generational indices were used, which allow for relations to be
pointed to while also automatically invalidating pointers to removed relations. (TODO: cite)

Each edge between agents and relations and vice versa also contains a small amount of
utility information that helps different world operations complete their tasks
efficiently, such as storing the index an agent occupies within a relation to identify
which edge of a directional it is on.

(TODO: visual)

\section{Type Checking} \label{sec:TypeChecking}

After the world is built, a scan is passed over it to ensure that expressions resolve
correctly to the type that is expected of them. For example, a predicate goal must
evaluate to a boolean value, and a delta goal must evaluate to a number.

The type checker also keeps track of guarantees, which are stored within any boolean
type that a function is determined to evaluate to. For example, the expression
$alice.feels.happy$ evaluates to a boolean type, but it also carries the guarantee that
if it is true, the $alice.feels.happy$ will exist. Then, since the right hand of an
$and$ expression is only evaluated if the left hand side is $true$, we can guarantee
in the right hand branch that $alice.feels.happy$ will exist and can have its fields
accessed.

(TODO: visual)

Type checking ended up being a very necessary part of the system. Every major change
or addition would typically result in a few type checker errors that had to be resolved.
By nature, these simulations are highly branching which makes them hard to thoroughly
test. Catching these errors during world build time reduces latent logic errors
significantly and offers useful feedback on how to achieve certain outcomes.

\section{The Simulation Layer} \label{sec:Transactions}

The world only stores data and simple utility functions for adding and modifying
relations, which allows for a simulation layer to be built on top of it that handles
common operations that are performed in more practical scenarios. For example,
the simulation function for getting the value of a relation's field constructs a
query from the sentence, ensures it is well-formed and refers to a relation type
and actors that exist, and then forwards that request to the world, propagating
and adding context to any errors it encounters along the way.

Simulation functions that modify the world, such as most of the practice outcomes\/effects,
require an additional transaction layer to be added to the pipeline. This layer records
all low-level mutations made to the world, and allows for snapshotting and rollbacks of
the world state to a specific point in time. By storing changes to the world state instead
of distinct copies of the world state, the world can be very quickly iterated over
to evaluate the scores of future states for the planning system. This is one large
performance enhancement over the original Praxis, which used copies for evaluating actions (TODO: cite),
and opens the door for planning further into the future.

(TODO: visual)

\section{Planning System} \label{sec:PlanningSystem}

The Praxsmith engine includes a built-in planner that uses the goals defined for each
actor to score their potential actions by desirability. For example, if an actor has a
goal with a weight of $5$ that counts all of the happy people in the world, actions that
make people happy will be scored with $5$ per person, including anyone who becomes unhappy
counting as negative score.

As stated previously, the original Praxis created a new world state evaluated for each
action and ran goal evaluations on each of those states to determine the score for
each action. This is an elegant solution from a functional standpoint, but it does
introduce the inefficiency of cloning the world every time an action is evaluated.

Praxsmith, on the other hand, keeps track of every change in the world state and
provides the ability to roll world state back through a transactional interface. For
every action, the current point in the change log is recorded, the action is evaluated
and scored, and then the world is rolled back to its initial state.

This efficiency allows for higher depths of evaluations to be calculated, which is
provided by a variable on the frontend when querying for agent actions. A depth of $1$
will not catch a string of actions making someone happy 5 moves out, but a depth of
$5$ will be able to think that far ahead. This allows agents to take more complex lines
of thought and practically limits them only by what goals they have set.

Introducing a ``Do Nothing'' no-op action exposes an issue with depth scoring, which is
that the chain of $Do Nothing \rightarrow Do Nothing \rightarrow Positive Action$
will score the same as just $Positive Action$ because they both end up in the same state,
even though one has multiple unnecessary steps. To avoid this problem, each change in score
decays when it is passed down the recursion change. In this case, the score delta from
the initial world state to the state after $Positive Action$ is processed might be $5$,
but after the two $Do Nothing$ actions, that score decays down to $3.2$, assuming a
decay rate of $0.8$. This allows the planner to prioritize quick positive results over ones
that may take several turns.

(TODO: figure)

One limitation of this system is that it only accounts for the actions your character
could take. This is primarily because the core Praxsmith engine has no notion of
actor turn order, so it wouldn't know whose actions to evaluate next. This is explained
more in the Limitations section (TODO: ref), but centralizing action order would allow
for characters to take the potential actions of others into account.
